% Generated draft: requires mathematical and source review.
\documentclass[11pt]{article}
\usepackage[margin=25mm]{geometry}
\usepackage{fontspec}
\setmainfont{Latin Modern Roman}
\usepackage{amsmath,amssymb,mathtools,unicode-math}
\setmathfont{Latin Modern Math}
\usepackage{enumitem,needspace,xurl,hyperref}
\setlength{\parindent}{0pt}
\setlength{\parskip}{5pt}
\setlength{\emergencystretch}{3em}
\newcommand{\N}{\mathbb{N}}
\newcommand{\Z}{\mathbb{Z}}
\newcommand{\Q}{\mathbb{Q}}
\newcommand{\R}{\mathbb{R}}
\newcommand{\C}{\mathbb{C}}
\newcommand{\F}{\mathbb{F}}
\newcommand{\A}{\mathbb{A}}
\newcommand{\PP}{\mathbb{P}}
\newcommand{\E}{\mathbb{E}}
\newcommand{\eps}{\varepsilon}
\newcommand{\dd}{\,\mathrm{d}}
\begin{document}
\section*{UnsolvedMath ID: 30003132}
\textbf{Source problem code:} catalogue entry
\subsection{Definitions and mathematical statement}
For notation, let \(n\geq 1\), let \(\Omega_1,\ldots,\Omega_n\) be nonempty finite sets, and let \(\mu_i\) be a probability measure on \(\Omega_i\). Write
\[
\Omega=\prod_{i=1}^n\Omega_i,
\qquad
\mathbb P=\bigotimes_{i=1}^n\mu_i.
\]
To interpret increasing events, coordinate orders must also be specified. Given partial orders \(\leq_i\) on the factors, use the coordinatewise order
\[
x\leq y\quad\Longleftrightarrow\quad x_i\leq_i y_i\text{ for every }i.
\]
An event \(D\subseteq\Omega\) is increasing, or an up-set, if \(x\in D\) and \(x\leq y\) imply \(y\in D\).

For increasing events \(A,B,C\subseteq\Omega\), define their seven nonempty intersection probabilities by
\[
\begin{aligned}
a&=\mathbb P(A),& b&=\mathbb P(B),& c&=\mathbb P(C),\\
ab_*&=\mathbb P(A\cap B),& ac_*&=\mathbb P(A\cap C),& bc_*&=\mathbb P(B\cap C),\\
t&=\mathbb P(A\cap B\cap C).
\end{aligned}
\]
Here \(ab_*\), \(ac_*\), and \(bc_*\) denote intersection probabilities, not products of marginal probabilities. Equivalently, to avoid this notation, define the realizable set for a specified ordered product space by
\[
\mathcal R(\Omega,\mathbb P)=
\left\{
\bigl(\mathbb P(A),\mathbb P(B),\mathbb P(C),
\mathbb P(A\cap B),\mathbb P(A\cap C),\mathbb P(B\cap C),
\mathbb P(A\cap B\cap C)\bigr):
A,B,C\subseteq\Omega\text{ are increasing}
\right\}.
\]
\textbf{Problem.} Characterize the realizable joint probabilities of three increasing events in a finite product probability space. In particular, determine whether, for every such space and every three increasing events \(A,B,C\),
\[
\begin{aligned}
&2\mathbb P(A\cap B\cap C)
-\mathbb P(A\cap B)\mathbb P(C)
-\mathbb P(A\cap C)\mathbb P(B)\\
&\qquad
-\mathbb P(B\cap C)\mathbb P(A)
+\mathbb P(A)\mathbb P(B)\mathbb P(C)
\geq 0.
\end{aligned}
\]
\textbf{Unresolved scope.} The supplied record does not specify the coordinate orders or the permitted class of ordered factors. The order-relative definitions above explain the terminology; they do not establish that arbitrary partially ordered factors are intended. It is also unspecified whether the requested characterization concerns \(\mathcal R(\Omega,\mathbb P)\) for each fixed space or the union of these sets as the permitted spaces vary. The universal quantifier in the inequality must therefore be read relative to the intended, presently unspecified class of ordered product spaces. The record imposes no strict-positivity condition on the factor probabilities and no exclusion of empty or full events.
\subsection{Short English statement}
Describe exactly which seven intersection probabilities can arise from three increasing events in a finite product probability space, and determine whether the displayed three-event inequality always holds. The permitted coordinate orders and the scope of the realizability characterization require clarification. The supplied record labels the problem open; this draft does not independently verify that status.
\subsection{Sources}
\begin{itemize}
\item Supplied local UnsolvedMath registry record, permanent ID \(30003132\), source problem code \texttt{catalogue entry}, title \emph{Three-Event Correlation Inequalities in Product Spaces}, in the Oberwolfach Reports Open Problems collection. The record gives proposed year 2016 and names no proposer.
\item Supplied report citation: \emph{Combinatorics and Probability (2017)}, Oberwolfach Reports, pp.~1189--1257, DOI \texttt{10.4171/owr/2016/22}. \url{https://doi.org/10.4171/owr/2016/22}. The supplied background describes Oberwolfach Reports as documenting mathematical workshops and their open-problem discussions. The citation's year 2017 is retained alongside the record's proposed year 2016 and the DOI containing 2016; their relationship has not been independently checked.
\item Sahi, \emph{Higher Correlation Inequalities}. Supplied literature source: \url{https://sites.math.rutgers.edu/~sahi/Reprints/08highercorrelation.pdf}.
\item Dubey, Sahi, and Wang, \emph{Putting All Eggs in One Basket}. Supplied literature source: \url{https://arxiv.org/abs/2403.15957}.
\item Literature assessment dated \texttt{2026-08-21}: the supplied curation labels the problem open and states that Sahi's three-event correlation inequality and a complete description of realizable joint probabilities for three up-sets remain open, while later work treats special forms without a general resolution. This is supplied dated curation, described in the record as literature triage derived from OWR curation, and has not been independently verified.
\item The supplied registry assigns statement status \texttt{corrected\_verified} and reports comparison with the official Oberwolfach report. That verification claim is preserved as source metadata; this draft does not independently confirm it.
\end{itemize}
\end{document}
